
\documentclass[twocolumn]{article}

\usepackage{preprint}
\usepackage[utf8]{inputenc}
\usepackage[T1]{fontenc}
\usepackage{booktabs}
\usepackage{amsmath}
\usepackage{microtype}
\usepackage{xcolor}
\usepackage{url}
\usepackage{hyperref}
\usepackage[numbers,square]{natbib}

% arXiv two-column template-style hyperlinks
\hypersetup{
  colorlinks=true,
  linkcolor=purple,
  urlcolor=blue,
  citecolor=cyan,
  anchorcolor=black
}

\bibliographystyle{unsrtnat}

\title{Reconstruction Attacks on PCEN Spectrograms: A Negative Result and a Small Edge-Audio Case Study}

\author{Alessandro Canonico\\Bocconi University, Milan}

\date{}

\begin{document}

\twocolumn[{\begin{@twocolumnfalse}
\maketitle

\begin{abstract}
Per-channel energy normalization (PCEN) is a useful adaptive preprocessing for
acoustic event detection, but it is sometimes treated as if the
many-to-one feature transformation also provides privacy. We test that assumption
directly. We implement an oracle and a blind inverse PCEN attacker, reconstruct
3 second speech windows with Griffin-Lim and an off-the-shelf HiFi-GAN vocoder,
and measure intelligibility, speaker leakage, PESQ, and STOI against a plain
log-mel baseline. On 200 LibriSpeech dev-clean clips, PCEN does not add a
measurable privacy advantage over plain mel: oracle PCEN plus Griffin-Lim gives
WER 0.101 [0.081, 0.122] versus 0.091 [0.072, 0.112] for mel, while HiFi-GAN
has identical WER 0.267 [0.229, 0.305] under PCEN and mel. A blind inverse
still gives WER 0.150 [0.123, 0.177] and speaker cosine 0.797 [0.789, 0.804].
Planted phrases, including password and card number strings, survive
the transform. We include the motivating edge-audio classifier as an audited
case study rather than a deployment claim and is not presented as deployment-ready: the model is compact
(1.83M parameters, 3.72 MB INT8 ONNX), but its base AUC is 0.804
[0.793, 0.815] and the only measured per-site result fails to statistically support the target false-positive and recall gates; we therefore report it as an engineering audit rather than a deployment result.
The main contribution is a
reproducible attack-and-claim-audit harness, together with a negative result
that prevents overclaiming PCEN as a privacy primitive.
\end{abstract}

\keywords{PCEN, audio privacy, reconstruction attacks, edge audio classification, claim auditing}
\vspace{0.5cm}
\end{@twocolumnfalse}}]

\section{Introduction}

PCEN was introduced as an adaptive time frequency preprocessing for acoustic event
detection and related audio tasks \cite{lostanlen2019pcen}. Its temporal
integration, adaptive gain control, and dynamic range compression, often improve
robustness over a static log-mel representation. The same transform is
many-to-one: phase is discarded, magnitudes are compressed, and an internal
running smoother is not normally exported.

Many-to-one is not the same as private. A defensible privacy statement needs a
threat model and an empirical attack evaluation. This repository originally
used PCEN as the input representation for a small edge classifier and treated
the transform too strongly in informal project framing. This paper records the
audit that corrected that claim.

The paper has two linked parts. First, it evaluates whether PCEN adds privacy
over a plain log-mel baseline against standard reconstruction attacks. Second,
it reports the compact classifier only as a case study in claim discipline: the
classifier is reproducible and efficient, but the data does not yet support
deployment or cross-site generalization claims.

\section{Threat Model}

The attacker observes the PCEN tile consumed by the classifier: 64 mel bands by
188 frames from a 3 second, 16 kHz mono window. The attacker knows the public
feature configuration, including the PCEN defaults used by \texttt{librosa.pcen}
with gain 0.98, bias 2.0, power 0.5, time constant 0.4 seconds, and
\(\epsilon=10^{-6}\), applied to a mel spectrogram scaled by \(2^{31}\).

The attacker does not observe the original waveform. In the realistic setting,
the attacker also does not observe the PCEN running mean state \(M\). We test
two inverse-PCEN settings:

\begin{itemize}
  \item \textbf{Oracle inverse}: exact closed-form inverse using the forward
  running mean \(M\). This is a strongest-attacker upper bound.
  \item \textbf{Blind inverse}: iterative inverse that estimates \(M\) from the
  PCEN tile itself. This is the realistic attacker.
\end{itemize}

We compare both settings against a mel-baseline condition that reconstructs
from the corresponding log-mel representation without PCEN. The PCEN-minus-mel
delta is the measured privacy contribution of PCEN.

\section{Attack Harness}

The evaluation is implemented as repository test code. The forward PCEN
implementation matches \texttt{librosa.pcen}; the oracle inverse round-trips
to floating-point precision, and the blind inverse converges in 10 to 15
iterations after peak normalization.

We test two waveform reconstruction attacks. The classical attack uses
Griffin-Lim phase recovery through \texttt{librosa.feature.inverse.mel\_to\_audio}
with 60 iterations \cite{griffinlim1984}. The neural attack uses the
SpeechBrain LibriTTS HiFi-GAN vocoder \cite{kong2020hifigan,speechbrain2021}.
Because the classifier frontend has 64 mel bands and the vocoder expects 80, we
use a public least-squares adapter from the 64-band basis to the 80-band basis.

The main corpus is 200 deterministic 3 second clips from LibriSpeech dev-clean
\cite{panayotov2015librispeech}. A separate hidden-phrase corpus contains five
TTS-generated clips with planted phrases such as passwords and
card number strings. Metrics are Whisper-tiny.en WER
\cite{radford2022whisper}, ECAPA-TDNN speaker cosine, PESQ wideband, and STOI.
All aggregate intervals are percentile bootstrap 95\% confidence intervals.

\section{Reconstruction Results}

Table~\ref{tab:privacy} gives the main LibriSpeech result. PCEN does not add a
measurable privacy advantage over plain mel. Under Griffin-Lim, oracle PCEN
has WER 0.101 [0.081, 0.122] versus mel WER 0.091 [0.072, 0.112]. Under
HiFi-GAN, PCEN and mel have identical WER 0.267 [0.229, 0.305]. The blind
attacker still reconstructs intelligible speech, with WER 0.150 [0.123,
0.177] and speaker cosine 0.797 [0.789, 0.804], far above the chance cosine
0.114 [0.095, 0.135].

\begin{table}[t]
\centering
\tiny
\setlength{\tabcolsep}{2pt}
\begin{tabular}{llcc}
\toprule
Recovery & Attack & WER & Speaker cosine \\
\midrule
pcen\_oracle & Griffin-Lim & 0.101 [0.081, 0.122] & 0.841 [0.835, 0.847] \\
pcen\_oracle & HiFi-GAN & 0.267 [0.229, 0.305] & 0.550 [0.539, 0.562] \\
pcen\_blind & Griffin-Lim & 0.150 [0.123, 0.177] & 0.797 [0.789, 0.804] \\
pcen\_blind & HiFi-GAN & 0.206 [0.175, 0.238] & 0.631 [0.621, 0.642] \\
mel\_baseline & Griffin-Lim & 0.091 [0.072, 0.112] & 0.838 [0.832, 0.845] \\
mel\_baseline & HiFi-GAN & 0.267 [0.229, 0.305] & 0.550 [0.539, 0.562] \\
\bottomrule
\end{tabular}
\caption{Reconstruction leakage on LibriSpeech dev-clean, n=200. Intervals are bootstrap 95\% CIs.}
\label{tab:privacy}
\end{table}

The hidden-phrase probes make the practical risk concrete. A phrase recognized
by Whisper on the original as ``credit card number is 1-2-3-4-5-6'' is
recovered verbatim across the reconstruction settings. A TTS pronunciation of
``fortepiano'' is recognized as ``4 to PM'' on the original and is often
recovered with the same text or a one-token spelling error. The leakage is not
only aggregate intelligibility; sensitive strings can survive the feature
pipeline.

\section{Classifier Case Study}

The motivating classifier is a CNN6-style encoder with squeeze-and-excitation,
a GRU temporal layer, and multi-scale pooling. It has 1,829,444 parameters, a
7.00 MB FP32 checkpoint, and a 3.72 MB static-INT8 ONNX export. Training uses
AudioSet, UrbanSound8K, ESC-50, YouTube-derived ambient and scream clips, and a
violence-audio dataset, with focal loss at \(\gamma=0.5\), cosine scheduling,
strong augmentation, and environmental noise injection.

The base classifier is not a validated deployment system. On the held-out test
split it reaches AUC-ROC 0.804 [0.793, 0.815], accuracy 70.34\%
[69.24, 71.41], threat recall 80.86\% [79.10, 82.50], and overall FPR
34.18\% [32.85, 35.54]. Hard negatives are the main failure mode: speech FPR
is 71.7\% [67.0, 76.0] and laughter FPR is 82.5\% [77.5, 86.6].

Per-site fine-tuning was explored on one metro site as an engineering audit,
not as an untouched final evaluation. The fine-tuning path used the base test
split as a proxy validation and threshold-selection set, and the metro
operating point was later selected from a checkpoint--threshold--aggregation
sweep. The selected audit point uses the \texttt{metro\_tweak3} recipe, the
pre-lever low-FPR threshold 0.601, and temporal-majority aggregation with
\(k=2\). It gives 0/19 ambient false positives and 45/57 threat detections,
but these are selected-point audit numbers rather than clean held-out
deployment estimates. The Wilson 95\% upper bound on false positives is
16.8\%, above the target 5\% gate. Threat recall is 78.9\% with Wilson 95\%
CI [66.7, 87.5], whose upper bound is below the 88\% target. Therefore
none of the deployment gates are statistically supported.

\section{Baselines and Efficiency}

The reproducible comparison against CNN14 is a footprint and latency
comparison, not a matched accuracy comparison. SafeCommute has 1.83M
parameters, while CNN14 has roughly 81.8M parameters, a 44.7x parameter ratio.
On the same Ryzen 7 7435HS reference machine, the 8-thread benchmark reports
median model-only latency 93.38 ms for CNN14 versus 6.09 ms for SafeCommute
INT8 ONNX, a 15.3x ratio. A separate 1-thread run reports 226.28 ms for CNN14
versus 23.66 ms for SafeCommute, a 9.6x ratio. End-to-end latency for
SafeCommute, including librosa mel and PCEN preprocessing, is much higher than
model-only latency: median 34.53 ms and p99 52.97 ms in the current 200-window
run. The paper therefore reports both model-only and end-to-end numbers.

\section{Data Availability and Reproducibility}

Data are available through public upstream datasets and local re-generation scripts where licensing permits. The public repository retains the code, paper source, and generated JSON reports needed to audit the PCEN reconstruction result. Raw audio, prepared tensors, model checkpoints, and generated waveform reconstructions are not included in the release snapshot:

\begin{itemize}
  \item The local YouTube-derived metro and scream WAVs cannot be redistributed
  as raw audio without a separate manifest and terms review.
  \item The local violence-audio folder has 2000 WAVs and is YouTube-derived, but video IDs, offsets, and
  redistribution terms are not captured in this repository.
  \item Prepared tensors inherit the restrictions of their source audio and
  should not be redistributed until all upstream licenses are cleared.
  \item Privacy evaluation corpora are reproducible from LibriSpeech dev-clean
  and generated hidden-phrase scripts, subject to the licenses of LibriSpeech
  and the TTS/vocoder models used to generate the samples.
  \item Generated or reconstructed privacy WAVs, if produced locally, are
  bounded audit artifacts. They are not raw training data, are not covered by
  the repository licenses, and remain subject to the upstream
  LibriSpeech/TTS/vocoder terms.
\end{itemize}

\section{Limitations}

This is a negative technical note, not a proof of a new privacy attack class.
The adjacent literature already treats mel-family features as reconstructable
attack surfaces rather than privacy mechanisms \cite{mestiri2024privacy,
safeear2024}. The PCEN result is useful primarily because it prevents a
specific overclaim and leaves a reusable audit harness.

The classifier study is also limited. It has one real per-site evaluation,
small held-out denominators, no cross-site confidence interval, and no matched
classifier quality comparison against large audio models. Its strongest
validated claims are compactness, reproducible metrics, and the failure of
previously stronger deployment claims under confidence intervals.

\section{Conclusion}

PCEN should not be described as a privacy primitive for edge audio classifiers.
Off-the-shelf reconstruction attacks recover intelligible speech and speaker
identity from PCEN tiles, and PCEN does not measurably improve over a plain
log-mel baseline in this evaluation. The associated classifier is compact and
fully auditable, but current detection results support only a research case
study with explicit limitations. The release should prioritize reproducibility
and claim hygiene over deployment framing.

\bibliography{references}

\end{document}
